\section{단위원, 영인수, 정역과 체}
\label{sec:ring-units-domains}

\begin{learninggoal}
이 절에서는 환의 원소를 곱셈 관점에서 분류한다. 단위원과 영인수를 찾고, 정역과 체의 차이를 설명하며, 유한 정역이 체가 되는 이유를 증명한다.
\end{learninggoal}

\subsection{단위원}

\begin{definition}[단위원]
환 $R$의 원소 $u$에 대해 어떤 $v\in R$가 존재하여 $uv=vu=1$이면 $u$를 \emph{단위원}(unit)이라 한다. 모든 단위원의 집합을 $R^\times$로 쓴다.
\end{definition}

단위원은 곱셈에 대해 군을 이룬다. 이 군은 환에 숨어 있는 중요한 군론적 자료이다.

\begin{proposition}
환 $R$에 대해 $R^\times$는 곱셈군이다.
\end{proposition}

\begin{proof}
$1$은 단위원이고, 두 단위원 $u,v$의 곱은 역원 $v^{-1}u^{-1}$을 가지므로 다시 단위원이다. 단위원의 역원 역시 단위원이다. 결합법칙은 환의 곱셈에서 물려받는다.
\end{proof}

\begin{example}
\begin{enumerate}[label=(\alph*)]
  \item $\Z^\times=\{1,-1\}$.
  \item 체 $K$에서는 $K^\times=K\setminus\{0\}$.
  \item $(\Z/n\Z)^\times$는 $n$과 서로소인 합동류들로 이루어진다.
  \item $R\times S$에서 $(r,s)$가 단위원일 필요충분조건은 $r\in R^\times$, $s\in S^\times$이다.
\end{enumerate}
\end{example}

\begin{proposition}[합동류의 단위원 판정]
$\overline a\in\Z/n\Z$가 단위원일 필요충분조건은 $\gcd(a,n)=1$이다.
\end{proposition}

\begin{proofidea}
$\overline a$가 역원을 가진다는 것은 어떤 $b$에 대해 $ab\equiv1\pmod n$이라는 뜻이다. 이는 어떤 정수 $k$에 대해 $ab+kn=1$이라는 베주 항등식과 같다.
\end{proofidea}

\begin{proof}
$\gcd(a,n)=1$이면 베주 항등식으로 $ax+ny=1$인 $x,y\in\Z$가 존재한다. 따라서 $\overline a\,\overline x=\overline1$이다.

반대로 $\overline a$가 단위원이면 $ab\equiv1\pmod n$인 $b$가 있으므로 $ab-1$은 $n$의 배수이다. 따라서 $ab+nk=1$인 $k\in\Z$가 존재하고, $a,n$의 모든 공약수는 $1$도 나누므로 $\gcd(a,n)=1$이다.
\end{proof}

\subsection{영인수와 멱영원소}

\begin{definition}[영인수]
가환환 $R$에서 $a\ne0$이고 어떤 $b\ne0$에 대해 $ab=0$이면 $a$를 \emph{영인수}(zero divisor)라 한다.
\end{definition}

\begin{definition}[멱영원소]
어떤 양의 정수 $m$에 대해 $a^m=0$이면 $a$를 \emph{멱영원소}(nilpotent element)라 한다.
\end{definition}

모든 영이 아닌 멱영원소는 영인수이다. 실제로 $m$을 $a^m=0$이 되는 최소 양의 정수라 하면 $a^{m-1}\ne0$이고 $a\cdot a^{m-1}=0$이다.

\begin{example}
$\Z/12\Z$에서 $\overline3\cdot\overline4=\overline0$이므로 $\overline3$과 $\overline4$는 영인수이다. 또한
\[
  \overline6^2=\overline0
\]
이므로 $\overline6$은 멱영원소이다. 반면 $\overline3$은 멱영원소가 아니다.
\end{example}

\begin{proposition}
환 $R$에서 $a$가 멱영원소이면 $1-a$는 단위원이다.
\end{proposition}

\begin{proof}
$a^m=0$라 하자. 유한 등비급수 공식으로
\[
  (1-a)(1+a+a^2+\cdots+a^{m-1})=1-a^m=1
\]
이다.
\end{proof}

\begin{pitfall}
영인수와 단위원은 서로 배타적이다. $u$가 단위원이고 $ux=0$이면 양변에 $u^{-1}$을 곱해 $x=0$을 얻는다. 그러나 영인수가 아닌 원소가 항상 단위원인 것은 아니다. 정수환에서 $2$는 영인수가 아니지만 단위원도 아니다.
\end{pitfall}

\subsection{정역과 체}

\begin{definition}[정역과 체]
영환이 아닌 가환환 $R$이 영이 아닌 영인수를 갖지 않으면 $R$을 \emph{정역}(integral domain)이라 한다. 영환이 아닌 가환환 $K$에서 모든 영이 아닌 원소가 단위원이면 $K$를 \emph{체}(field)라 한다.
\end{definition}

체는 정역이지만 역은 일반적으로 거짓이다.

\begin{proposition}[정역에서의 소거법칙]
$R$이 정역이고 $a\ne0$일 때
\[
  ab=ac\quad\Longrightarrow\quad b=c
\]
이다.
\end{proposition}

\begin{proof}
$ab=ac$이면 $a(b-c)=0$이다. $a\ne0$이고 정역에는 영이 아닌 영인수가 없으므로 $b-c=0$이다.
\end{proof}

\begin{theorem}[유한 정역은 체]
유한 정역 $R$은 체이다.
\end{theorem}

\begin{proofidea}
영이 아닌 $a\in R$에 대해 곱셈 함수 $m_a(x)=ax$를 생각한다. 정역의 소거법칙으로 이 함수는 단사이고, 유한집합의 단사 자기함수는 전사이다. 따라서 $1$이 상에 들어가 $ax=1$인 $x$가 존재한다.
\end{proofidea}

\begin{proof}
$a\in R\setminus\{0\}$를 고정하자. 함수
\[
  m_a:R\to R,\qquad x\mapsto ax
\]
는 $ax=ay$이면 소거법칙에 의해 $x=y$이므로 단사이다. $R$이 유한이므로 $m_a$는 전사이다. 따라서 어떤 $b\in R$에 대해 $ab=1$이다. 모든 영이 아닌 원소가 단위원이므로 $R$은 체이다.
\end{proof}

\begin{corollary}
양의 정수 $n$에 대해 다음은 동치이다.
\begin{enumerate}[label=(\roman*)]
  \item $n$은 소수이다.
  \item $\Z/n\Z$는 정역이다.
  \item $\Z/n\Z$는 체이다.
\end{enumerate}
\end{corollary}

\begin{proof}
$n$이 소수이면 유클리드 보조정리에 의해 $ab\equiv0\pmod n$에서 $a\equiv0$ 또는 $b\equiv0$이므로 정역이다. 유한 정역은 체이다. 반대로 $n=rs$가 $1<r,s<n$인 합성수이면 $\overline r\ne0$, $\overline s\ne0$이지만 $\overline r\,\overline s=0$이므로 정역이 아니다.
\end{proof}

\subsection{구조의 포함 관계}

가환환의 중요한 계층은 다음과 같다.
\[
  \text{체}\Longrightarrow\text{정역}\Longrightarrow\text{가환환}.
\]
각 역방향은 일반적으로 성립하지 않는다.
\begin{itemize}
  \item $\Z$는 정역이지만 체가 아니다.
  \item $\Z/6\Z$는 가환환이지만 정역이 아니다.
\end{itemize}

\begin{checkpoint}
다음 환에서 단위원과 영인수를 구하라.
\[
  \Z/8\Z,\qquad \Z/10\Z,\qquad \Z/11\Z.
\]
그 결과만 보고 어느 환이 체인지 판단하라.
\end{checkpoint}
